\chapter{Visualisation des donnees : lecture et interpretation des graphiques}
\label{chap:visualisation}

Le tableau de bord de la plateforme (page \url{/dashboard}, alimentee par
\texttt{Dashboard\-Metrics\-Service} et rafraichie par le flux SSE decrit au
\cref{chap:realisation}) restitue la posture de risque du parc a travers cinq
graphiques complementaires, chacun repondant a une question de pilotage differente.
Ce chapitre reproduit chacun d'eux avec des \textbf{donnees illustratives} (le
tableau de bord reel affiche les donnees issues des imports Qualys/CMDB) et explique
comment le lire.

\section{Repartition par classe de risque}

\textbf{Question a laquelle il repond :} \emph{quelle proportion du parc se trouve
dans chaque classe de risque ?}

\begin{figure}[H]
\centering
\begin{tikzpicture}
\def\R{2.3}
% BLACK 5, RED 12, ORANGE 30, YELLOW 45, GREEN 120 -> total 212
\draw[fill=riskblack,draw=white]  (0,0) -- (90:\R)     arc (90:81.51:\R)    -- cycle;
\draw[fill=riskred,draw=white]    (0,0) -- (81.51:\R)  arc (81.51:61.13:\R) -- cycle;
\draw[fill=riskorange,draw=white] (0,0) -- (61.13:\R)  arc (61.13:10.19:\R) -- cycle;
\draw[fill=riskyellow,draw=white] (0,0) -- (10.19:\R)  arc (10.19:-66.23:\R)-- cycle;
\draw[fill=riskgreen,draw=white]  (0,0) -- (-66.23:\R) arc (-66.23:-270:\R)-- cycle;
\draw[fill=white] (0,0) circle (1.1); % anneau (doughnut)

% legende
\foreach \color/\label/\y in {riskblack/BLACK (5)/2.1, riskred/RED (12)/1.6,
   riskorange/ORANGE (30)/1.1, riskyellow/YELLOW (45)/0.6, riskgreen/GREEN (120)/0.1} {
  \node[draw,fill=\color,minimum width=0.3cm,minimum height=0.3cm] at (4.3,\y) {};
  \node[right=0.2cm,font=\scriptsize] at (4.6,\y) {\label};
}
\end{tikzpicture}
\caption{Repartition des actifs par classe de risque (diagramme en anneau, donnees illustratives)}
\label{fig:chart-class}
\end{figure}

\textbf{Comment le lire.} Chaque secteur represente une classe de risque
(\cref{tab:sla}) ; sa taille angulaire est proportionnelle au nombre d'actifs qui y
appartiennent, et sa couleur suit le code couleur officiel de la matrice SLA (noir,
rouge, orange, jaune, vert). Un coup d'oeil suffit pour evaluer la sante globale du
parc : un anneau majoritairement vert avec des secteurs noir/rouge fins indique un
parc globalement maitrise avec un nombre restreint d'urgences a traiter -- c'est la
lecture attendue sur l'exemple ci-dessus, ou les classes BLACK et RED ne
representent ensemble que $17/212\approx 8\%$ du parc.

\section{Repartition par equipe proprietaire}

\textbf{Question a laquelle il repond :} \emph{quelle equipe porte le plus de risque
critique a traiter ?}

\begin{figure}[H]
\centering
\begin{tikzpicture}
\begin{axis}[
  xbar stacked, width=13cm, height=5cm,
  xlabel={Nombre d'actifs}, xmin=0,
  symbolic y coords={Equipe C,Equipe B,Equipe A},
  ytick=data, bar width=8mm,
  legend style={at={(0.5,-0.28)},anchor=north,legend columns=5,font=\scriptsize},
  legend cell align=left
]
\addplot [fill=riskgreen]  coordinates {(40,Equipe A) (55,Equipe B) (25,Equipe C)};
\addplot [fill=riskyellow] coordinates {(10,Equipe A) (15,Equipe B) (20,Equipe C)};
\addplot [fill=riskorange] coordinates {(5,Equipe A) (10,Equipe B) (15,Equipe C)};
\addplot [fill=riskred]    coordinates {(2,Equipe A) (4,Equipe B) (6,Equipe C)};
\addplot [fill=riskblack]  coordinates {(0,Equipe A) (1,Equipe B) (4,Equipe C)};
\legend{Green,Yellow,Orange,Red,Black}
\end{axis}
\end{tikzpicture}
\caption{Repartition des actifs par equipe proprietaire (barres horizontales empilees, donnees illustratives)}
\label{fig:chart-owner}
\end{figure}

\textbf{Comment le lire.} Chaque ligne horizontale correspond au proprietaire CMDB
(\texttt{cmdbOwner}) d'un ensemble d'actifs ; sa longueur totale est le nombre
d'actifs qu'elle possede, et chaque segment colore est le nombre d'actifs de cette
equipe dans une classe de risque donnee. Contrairement au graphique en anneau
(vision globale), ce graphique permet de \textbf{cibler l'action} : sur l'exemple,
l'Equipe C possede moins d'actifs au total que l'Equipe B, mais une proportion bien
plus grande en classes Orange/Red/Black -- c'est elle qui devrait etre prioritairement
sollicitee par le comite de securite.

\section{Flux Environnement \texorpdfstring{$\to$}{->} Role \texorpdfstring{$\to$}{->} Classe de risque}

\textbf{Question a laquelle il repond :} \emph{d'ou vient le risque, structurellement
-- quels environnements et quels roles techniques concentrent les actifs les plus
critiques ?}

\begin{figure}[H]
\centering
\begin{tikzpicture}[font=\scriptsize]
% Colonnes : Environnement -> Role -> Classe
\node[draw,fill=dxcblue!15,minimum width=2.4cm,minimum height=0.7cm] (prod) at (0,1.3) {PRODUCTION};
\node[draw,fill=dxcblue!15,minimum width=2.4cm,minimum height=0.7cm] (nonprod) at (0,-1.1) {STAGING / TEST / DEV};

\node[draw,fill=dxcgray!12,minimum width=2.6cm,minimum height=0.7cm] (db) at (5.2,2.1) {DATABASE\_SERVER};
\node[draw,fill=dxcgray!12,minimum width=2.6cm,minimum height=0.7cm] (vm) at (5.2,0.3) {VS\_LINUX/WINDOWS};
\node[draw,fill=dxcgray!12,minimum width=2.6cm,minimum height=0.7cm] (net) at (5.2,-1.6) {IP\_ROUTER/SWITCH};

\node[draw,fill=riskred!70,minimum width=1.6cm,minimum height=0.6cm] (red) at (10.4,2.5) {RED};
\node[draw,fill=riskorange!70,minimum width=1.6cm,minimum height=0.6cm] (org) at (10.4,1.3) {ORANGE};
\node[draw,fill=riskyellow!70,minimum width=1.6cm,minimum height=0.6cm] (yel) at (10.4,0.1) {YELLOW};
\node[draw,fill=riskgreen!70,minimum width=1.6cm,minimum height=0.6cm] (grn) at (10.4,-1.3) {GREEN};

% Flux Env -> Role (epaisseur ~ volume)
\draw[riskred,line width=3pt,opacity=0.45] (prod.east) to[out=0,in=180] (db.west);
\draw[dxcblue,line width=6pt,opacity=0.35] (prod.east) to[out=0,in=180] (vm.west);
\draw[dxcgray,line width=2pt,opacity=0.4] (prod.east) to[out=0,in=180] (net.west);
\draw[dxcblue,line width=8pt,opacity=0.3] (nonprod.east) to[out=0,in=180] (vm.west);
\draw[dxcgray,line width=2pt,opacity=0.4] (nonprod.east) to[out=0,in=180] (net.west);

% Flux Role -> Classe
\draw[riskred,line width=3pt,opacity=0.5] (db.east) to[out=0,in=180] (red.west);
\draw[riskorange,line width=2.5pt,opacity=0.5] (db.east) to[out=0,in=180] (org.west);
\draw[riskorange,line width=3pt,opacity=0.5] (vm.east) to[out=0,in=180] (org.west);
\draw[riskyellow,line width=5pt,opacity=0.5] (vm.east) to[out=0,in=180] (yel.west);
\draw[riskgreen,line width=4pt,opacity=0.5] (vm.east) to[out=0,in=180] (grn.west);
\draw[riskyellow,line width=2pt,opacity=0.5] (net.east) to[out=0,in=180] (yel.west);
\draw[riskgreen,line width=5pt,opacity=0.5] (net.east) to[out=0,in=180] (grn.west);

\node[above=0.1cm of prod, font=\scriptsize\bfseries] {\phantom{x}};
\node[font=\scriptsize\bfseries] at (0,2.5) {Environnement};
\node[font=\scriptsize\bfseries] at (5.2,3.2) {Role technique};
\node[font=\scriptsize\bfseries] at (10.4,3.2) {Classe de risque};
\end{tikzpicture}
\caption{Diagramme de flux (Sankey) Environnement $\to$ Role $\to$ Classe de risque (donnees illustratives)}
\label{fig:chart-sankey}
\end{figure}

\textbf{Comment le lire.} Chaque bande horizontale relie deux categories ; sa
\textbf{largeur (l'epaisseur du trait) est proportionnelle au nombre d'actifs}
empruntant ce chemin, et sa couleur reprend soit l'environnement d'origine, soit --
pour le second niveau -- la classe de risque de destination (c'est exactement la
convention utilisee par \texttt{DashboardMetricsService}, qui construit ce graphique
en deux temps : \texttt{Env}$\to$\texttt{Role} d'abord, \texttt{Role}$\to$\texttt{RiskClass}
ensuite). On lit ce graphique de gauche a droite en suivant l'epaisseur des flux : sur
l'exemple, les bases de donnees de production concentrent un flux epais vers RED,
signalant qu'il s'agit du couple (role, environnement) le plus critique du parc --
une information que ni le graphique en anneau ni le graphique par equipe ne
permettent de voir directement, puisqu'aucun des deux ne croise simultanement
l'environnement et le role technique.

\section{Historique de la repartition du parc}

\textbf{Question a laquelle il repond :} \emph{la posture de risque du parc
s'ameliore-t-elle ou se degrade-t-elle au fil des campagnes de scan successives ?}

\begin{figure}[H]
\centering
\begin{tikzpicture}
\begin{axis}[
  width=13cm, height=6cm,
  xlabel={Campagne de scan (n)}, ylabel={Nombre d'actifs},
  xtick={1,2,3,4,5}, xticklabels={n-4,n-3,n-2,n-1,n},
  ymin=0, stack plots=y, area style,
  legend style={at={(0.5,-0.22)},anchor=north,legend columns=5,font=\scriptsize}
]
\addplot [fill=riskblack, draw=riskblack]  coordinates {(1,4)(2,5)(3,6)(4,5)(5,5)} \closedcycle;
\addplot [fill=riskred, draw=riskred]      coordinates {(1,18)(2,16)(3,15)(4,13)(5,12)} \closedcycle;
\addplot [fill=riskorange, draw=riskorange]coordinates {(1,40)(2,38)(3,35)(4,32)(5,30)} \closedcycle;
\addplot [fill=riskyellow, draw=riskyellow]coordinates {(1,50)(2,48)(3,47)(4,46)(5,45)} \closedcycle;
\addplot [fill=riskgreen, draw=riskgreen]  coordinates {(1,90)(2,100)(3,108)(4,115)(5,120)} \closedcycle;
\legend{Black,Red,Orange,Yellow,Green}
\end{axis}
\end{tikzpicture}
\caption{Historique empile de la repartition par classe (donnees illustratives, une serie par instantane \texttt{RiskSnapshot})}
\label{fig:chart-history}
\end{figure}

\textbf{Comment le lire.} Chaque bande coloree empilee represente l'effectif d'une
classe de risque a un instant donne (un \texttt{RiskSnapshot} enregistre a chaque
recalcul, \cref{chap:conception}) ; la hauteur totale de la pile est la taille du
parc scanne a ce moment-la. Ce n'est pas la hauteur d'une bande prise isolement qui
importe, mais son \textbf{evolution relative dans le temps} : sur l'exemple, la bande
verte (GREEN) s'epaissit progressivement tandis que les bandes orange et rouge
s'amincissent -- signe d'une tendance favorable, meme si le nombre d'actifs en
classe Black reste globalement stable, ce qui justifierait une attention particuliere
portee a ce sous-ensemble specifique plutot qu'a la tendance globale.

\section{Conformite SLA (age vs. score de risque)}

\textbf{Question a laquelle il repond :} \emph{quels actifs sont en violation de
leur delai de remediation, et de combien ?}

\begin{figure}[H]
\centering
\begin{tikzpicture}
\begin{axis}[
  width=13cm, height=7cm,
  xlabel={Asset Risk Score (ARS)}, ylabel={Anciennete de la vulnerabilite ouverte (jours)},
  xmin=0, xmax=1000, ymin=0, ymax=100,
  axis on top
]
% Bandes de fond par classe (zones SLA)
\fill[riskgreen,opacity=0.12]  (axis cs:0,0)   rectangle (axis cs:200,100);
\fill[riskyellow,opacity=0.12] (axis cs:200,0) rectangle (axis cs:400,100);
\fill[riskorange,opacity=0.12] (axis cs:400,0) rectangle (axis cs:600,100);
\fill[riskred,opacity=0.12]    (axis cs:600,0) rectangle (axis cs:700,100);
\fill[riskblack,opacity=0.12]  (axis cs:700,0) rectangle (axis cs:1000,100);

% Ligne en escalier : delai SLA exige par zone (90/30/15/7/2 jours)
\addplot[thick,dashed,dxcgray,const plot] coordinates {(0,90)(200,30)(400,15)(600,7)(700,2)(1000,2)};

% Points : actifs (vert = ON_TRACK, rouge = VIOLATED)
\addplot[only marks,mark=*,mark options={fill=riskgreen,draw=black,opacity=0.85}]
  coordinates {(90,20)(150,15)(250,10)(320,25)(480,8)(550,10)};
\addplot[only marks,mark=*,mark options={fill=riskred!80!black,draw=black,opacity=0.9}]
  coordinates {(260,45)(430,20)(610,12)(650,18)(720,6)(760,9)};

\legend{Limite SLA (escalier), Actifs conformes (ON\_TRACK), Actifs en violation (VIOLATED)}
\end{axis}
\end{tikzpicture}
\caption{Conformite SLA : anciennete vs. ARS, avec zones de classe en fond (donnees illustratives)}
\label{fig:chart-sla}
\end{figure}

\textbf{Comment le lire.} Chaque point est un actif ; sa position horizontale est son
ARS (\cref{eq:ars}) et sa position verticale, l'age en jours de sa vulnerabilite
ouverte la plus ancienne. Les bandes de fond delimitent les cinq classes de risque
(\cref{tab:sla}) et la ligne en escalier grise trace le delai SLA maximal autorise
pour la classe correspondant a chaque plage d'ARS. \textbf{La regle de lecture est
geometrique} : un point situe \emph{au-dessus} de la ligne en escalier, dans sa
propre bande de couleur, est en violation de son SLA -- c'est exactement le critere
programme dans \texttt{applyRiskScore} (\cref{eq:sla-status}), rendu visuellement
immediat. Sur l'exemple, le point a $(ARS=260, \text{age}=45\text{j})$ tombe dans la
zone Yellow (SLA de $30$ jours) alors qu'il affiche $45$ jours d'anciennete : il est
donc bien au-dessus de la marche de l'escalier a cet endroit, coherent avec sa
couleur rouge (violation). A l'inverse, un point situe sous la ligne, dans sa bande,
est encore \emph{dans les temps}. Ce graphique est celui qui traduit le plus
directement la gouvernance SLA du \cref{chap:modelisation} en decision operationnelle
: il permet d'identifier en un coup d'oeil les actifs qui necessitent une escalade
immediate.
